\section{Localized statistics and the signal of one pair}
\label{sec:localized}

Bulk dilution suggests localizing in height, but localization changes
the arithmetic estimate required. We develop two finite observables:
a nonnegative coincidence count, and an analytic quadratic form with
transverse amplification. Their lower signals are exact. The upper
estimates needed for a contradiction remain separate, unproved inputs.

\subsection{Height windows, resolution and normalization}
Fix \(T>1\), a center \(v\in[T,2T]\), and \(0<h<T/2\). Let
\(L=(\log T)/(2\pi)\) and choose \(w\in C_c^\infty(\R)\) with
\(0\le w\le1\), support in \([-1,1]\), and \(w=1\) on
\([-1/2,1/2]\). For all nontrivial zeta zeros
\(\rho_j=\sigma_j+\ii t_j\), counted with multiplicity, put
\begin{equation}\label{eq:local-window}
 w_j(v)=w\!\left(\frac{t_j-v}{h}\right),\qquad
 q=hL,\qquad C_w(b)=\int_\R w(r)w(r-b)\dd r,\qquad M=qC_w(0).
\end{equation}
Every sum below is finite because of the window. The quantity \(M\)
is a declared density normalization, not the zero count of every window.
The width \(h\) and the resolution of differences within a window
are distinct choices.

A final pair \(x_f\pm\ii d\) of \(H_0\), with \(d>0\),
corresponds to the two zeta zeros
\begin{equation}\label{eq:localized-pair}
 \rho_\pm=\frac12\pm\frac d2+\ii\theta,\qquad
 \theta=x_f/2,\qquad 0<d<1.
\end{equation}
They have exactly the same ordinate. Theorem~\ref{thm:tube}, if its
hypotheses hold to time zero, bounds both \(d\) and \(\theta\).
The common weight is \(w_0(v)=w((\theta-v)/h)\); it equals one on
the center interval \(J_\theta=[\theta-h/2,\theta+h/2]\).

\subsection{A positive coincidence signal}
Let \(\varphi\) be an even, nonnegative Schwartz function with
\(\varphi(0)=1\). For resolution \(\varepsilon>0\), define the
unnormalized ordered-pair statistic
\begin{equation}\label{eq:local-S}
 S_{v,h,\varepsilon}
 =\sum_{j\ne k}w_j(v)w_k(v)
       \varphi\!\left(\frac{L(t_j-t_k)}{\varepsilon}\right).
\end{equation}
The condition \(j\ne k\) removes a copy's self-pair, but retains
distinct copies at the same ordinate, including multiplicity.

\begin{proposition}[The local coincidence atom]\label{prop:local-atom}
If \(m_t\) is the total multiplicity at ordinate \(t\), then
\begin{equation}\label{eq:local-D}
 D_{v,h}:=\lim_{\varepsilon\downarrow0}S_{v,h,\varepsilon}
 =\sum_t w\!\left(\frac{t-v}{h}\right)^2m_t(m_t-1),
 \qquad S_{v,h,\varepsilon}\ge D_{v,h}.
\end{equation}
The two selected copies in \eqref{eq:localized-pair} contribute exactly
\(2w_0(v)^2\) to the coincidence part, and therefore
\(S_{v,h,\varepsilon}\ge2\) on \(J_\theta\).
\end{proposition}
\begin{proof}
In the finite sum, each term with \(t_j\ne t_k\) tends to zero,
whereas each equal-ordinate term has test value one. There are
\(m_t(m_t-1)\) ordered distinct pairs at ordinate \(t\).
All remaining summands are nonnegative, proving the inequality.
The selected pair contributes its two orders; any additional copies
at the same ordinate can only increase the statistic.
\end{proof}

Dividing by \(M\) gives a pair contribution \(2w_0^2/M\).
It remains of order one for \(hL\asymp1\), whereas it is diluted
when \(hL\to\infty\). Localization therefore evades the specific
bulk sensitivity limitation of Proposition~\ref{prop:finite-defects}.
It does not automatically evade averaging over centers. The exact
center integral of this contribution is
\begin{equation}\label{eq:atom-center-moment}
 \int_\R(2w_0(v)^2)^p\dd v
      =2^p h\int_\R w(r)^{2p}\dd r\qquad(p>0).
\end{equation}
If the profile is contained in \([T,2T]\), its normalized center
average thus has size \(h/T\), for every fixed \(p\).

The atom records multiplicity as well as off-line pairs. If
\(D_{v,h}=0\) for a covering family of windows, every ordinate
contains at most one simple zero; reflection then puts that zero on
the critical line. Conversely, simple critical-line zeros give
\(D_{v,h}=0\). This all-window assertion is RH together with
simplicity, and is not an independently established estimate.
Moreover, the coincidence signal does not grow with \(d\).

\subsection{The exact reference-process comparison}
To make the scale of a possible arithmetic estimate explicit, let
\(R_2(r)=1-(\sin\pi r/(\pi r))^2\), the unit-density sine-process
pair density, and put
\begin{equation}\label{eq:I-epsilon}
 I_\varepsilon=\int_\R\varphi(r/\varepsilon)R_2(r)\dd r
 =C_\varphi\varepsilon^3+O_\varphi(\varepsilon^5),\qquad
 C_\varphi=\frac{\pi^2}{3}\int_\R r^2\varphi(r)\dd r>0.
\end{equation}
Indeed \(R_2(r)=\pi^2r^2/3+O(r^4)\), with a global remainder
bound after enlarging the constant; the change of variable
\(r=\varepsilon u\) and the fourth Schwartz moment prove the expansion.

\begin{proposition}[Finite-window sine comparison]\label{prop:sine-window}
For the sine process, the expected value of \eqref{eq:local-S}, with
window width \(q\) in unit-density coordinates, is
\begin{equation}\label{eq:sine-window}
 \mathcal M_{q,\varepsilon}
 =q\int_\R C_w(r/q)\varphi(r/\varepsilon)R_2(r)\dd r
 \le M I_\varepsilon.
\end{equation}
More precisely, for \(q>0\) and \(0<\varepsilon\le1\),
\begin{equation}\label{eq:sine-window-error}
 0\le M I_\varepsilon-\mathcal M_{q,\varepsilon}
       \ll_{w,\varphi}\varepsilon^5/q.
\end{equation}
\end{proposition}
\begin{proof}
The expectation is the double integral of
\(w(a/q)w(b/q)\varphi((a-b)/\varepsilon)R_2(a-b)\).
Set \(r=a-b\) and integrate in \(a\) to obtain
\eqref{eq:sine-window}. Cauchy--Schwarz gives
\(C_w(b)\le C_w(0)\). Translation invariance and the fundamental
theorem of calculus give the sharper identity and bound
\[
 C_w(0)-C_w(b)=\tfrac12\|w-w(\cdot-b)\|_2^2
                       \le\tfrac12 b^2\|w'\|_2^2.
\]
Insert this bound in the integral and use \(0\le R_2(r)\ll r^2\).
The remaining fourth moment scales as \(\varepsilon^5\), proving
\eqref{eq:sine-window-error}.
\end{proof}

This proposition computes a reference-process main term; it asserts
no local zeta asymptotic. If one could independently prove
\begin{equation}\label{eq:coincidence-target}
 S_{v,h,\varepsilon}\le M I_\varepsilon+E(T,h,\varepsilon)
\end{equation}
uniformly at the relevant centers, a pair on the plateau would be
excluded only when the entire right-hand side is less than two.
For large \(M\), the main term suggests
\(\varepsilon\lesssim M^{-1/3}\), with a sufficiently small constant,
but an \emph{absolute} remainder below the residual one-pair margin
is still needed. A normalized \(o(1)\) error yields only \(o(M)\)
before normalization and is insufficient in general.
Also \(\widehat{\varphi(\cdot/\varepsilon)}(\xi)
=\varepsilon\widehat\varphi(\varepsilon\xi)\): narrowing the
test expands its Fourier scale by \(\varepsilon^{-1}\).
Theorem~\ref{thm:BGST} supplies neither this extended frequency range
nor the required moving-window estimate; its complex-difference
observable also differs from \(S\). Such an unmarked bound may
also exclude very close real pairs, a stronger demand than detecting
only nonreal zeros.

\subsection{An analytic signal with transverse growth}
For frequency \(a\ge0\), retain the kernel \(W\) of
\eqref{eq:F-complex} and define
\begin{equation}\label{eq:localized-E}
 E_{v,h}(a)=\sum_{j,k}w_j(v)w_k(v)
       \ee^{a(\rho_j-\rho_k)}W(\rho_j-\rho_k).
\end{equation}
This includes self terms. It is real and nonnegative, although its
individual summands need not be.

\begin{proposition}[Weighted Gram identity and background cost]
\label{prop:weighted-Gram}
For the finite reflected zero multiset in a height window,
\begin{equation}\label{eq:weighted-Gram}
 E_{v,h}(a)=\int_\R\ee^{-2|r|}
 \left|\sum_jw_j(v)\ee^{(\rho_j-1/2)(a+r)}\right|^2\dd r.
\end{equation}
Separate the two selected copies \eqref{eq:localized-pair} from the
remaining sum in this Hilbert norm. Their squared norm is
\begin{equation}\label{eq:pair-P}
 P_v=2w_0(v)^2[1+W(d)\cosh(ad)].
\end{equation}
If \(B_v\) is the squared norm of the remaining sum, then
\begin{equation}\label{eq:background-lower}
 E_{v,h}(a)\ge(\sqrt{P_v}-\sqrt{B_v})_+^2,
 \qquad
 P_v\le\bigl(\sqrt{E_{v,h}(a)}+\sqrt{B_v}\bigr)^2.
\end{equation}
\end{proposition}
\begin{proof}
The weights depend only on ordinate and are therefore equal on
\(\rho\mapsto1-\overline\rho\) partners. Relabel the second index
as in Proposition~\ref{prop:complex-Gram} and use
\(W(u)=\int_\R\ee^{-2|r|+ur}\dd r\). This proves
\eqref{eq:weighted-Gram}; all sums are finite and
\(|\Re u|<1\) ensures convergence. For the pair there are two
self terms, each \(w_0^2\), and two cross terms
\(w_0^2 W(d)\ee^{\pm ad}\), giving \eqref{eq:pair-P}.
Finally apply the triangle and reverse triangle inequalities to the
pair and background vectors in \(L^2(\ee^{-2|r|}\dd r)\).
\end{proof}

Thus positivity of \(E\) does not imply \(E\ge P_v\).
If independent bounds give \(E\le\mathcal U\) and
\(B_v\le\mathcal V\), the necessary comparison is
\(P_v\le(\sqrt{\mathcal U}+\sqrt{\mathcal V})^2\).
Both bounds are required. When all background zeros are on the critical
line, the triangle inequality gives
\(\sqrt{B_v}\le\sum_{\rm background}w_j\), because each exponential
has unit norm. This is an illustration under an additional hypothesis,
not an assumption in the proposed unconditional strategy.

For \(a=\alpha\log T\) and fixed \(d>0\), the pair norm on its
plateau grows like \(T^{\alpha d}\). If \(d=O(1/\log T)\) and
\(\alpha\) is bounded, this amplification stays bounded. A simpler
nonnegative transverse observable is
\begin{equation}\label{eq:local-transverse}
 A_{v,h}=\sum_jw_j(v)^2 L^2(\sigma_j-\tfrac12)^2.
\end{equation}
The pair contributes \(L^2d^2w_0^2/2\). Near a collision,
\eqref{eq:postcollision-jet} therefore identifies the dimensionless
time scale \(sL^2\). Vanishing of \(A\) in every covering window
is exactly RH; the definition provides no independent upper bound.

\subsection{Combining transport, amplitude and center width}
The preceding results can be assembled without hiding their missing
arithmetic premise. Let the controlled tube reach time zero and choose
\([T,2T]\) to contain its possible final ordinate. Let
\(J\subset J_\theta\cap[T,2T]\) be an interval of length \(\ell_J>0\).
On it, \(w_0=1\), and monotonicity of \(W(d)\) and \(\cosh(ad)\)
for \(0\le d<1\) gives
\begin{equation}\label{eq:P-star}
 P_v\ge P_*:=2[1+W(d_*)\cosh(ad_*)],
\end{equation}
provided the tube lower bound \(d_*\) is less than one. If
\(d_*\ge1\), the tube hypotheses already contradict the canonical
open strip.

\begin{corollary}[Conditional localized comparisons]
\label{cor:localized-contract}
Under these transport and plateau hypotheses:
\begin{enumerate}[label=(\roman*)]
\item A bound \eqref{eq:coincidence-target} at any center in \(J\)
with right-hand side below two contradicts the pair.
\item If \(B_v\le\mathcal V\) on \(J\), put
\(A_*=(\sqrt{P_*}-\sqrt{\mathcal V})_+^2\).
Then \(E_{v,h}(a)\ge A_*\) there. A pointwise upper bound
\(E_{v,h}(a)\le\mathcal U<A_*\) at a center in \(J\) is impossible.
\item For any \(p>0\), the same background bound implies
\begin{equation}\label{eq:moment-lower}
 \frac1T\int_T^{2T} E_{v,h}(a)^p\dd v
                  \ge\frac{\ell_J}{T}A_*^p.
\end{equation}
An independent moment upper bound strictly below this allowance
also contradicts the pair. Similarly,
\(T^{-1}\int_T^{2T}S_{v,h,\varepsilon}^p\dd v
\ge(\ell_J/T)2^p\).
\end{enumerate}
\end{corollary}
\begin{proof}
Apply Propositions~\ref{prop:local-atom} and~\ref{prop:weighted-Gram}
on the plateau, and integrate the resulting nonnegative lower bounds
over \(J\).
\end{proof}
This corollary is a conditional assembly of proved finite identities,
not an established arithmetic bound. If \(\ell_J\asymp h\) and a
signal comparable to \(T^{\alpha d}\) survives the background, the
moment lower scale is \((h/T)T^{p\alpha d}\). Transverse growth can
therefore overcome center dilution in principle. The background estimate
and an unconditional moment theorem at the same window and frequency
are the substantive missing inputs. A fixed-height exception does not
by itself supply a sequence of exceptions tending to infinite height;
an asymptotic argument must also address its finite-height range.

\subsection{What the available local estimates pay}
The relevant published statements have different quantifiers.
Maples--Rodgers' unconditional mesoscopic central limit theorem uses
a uniformly random center in \([T,2T]\), compact bounded-variation
tests with diverging variance, and \(n(T)\to\infty\) with
\(n(T)=o(\log T)\) \cite[Theorem~1.2]{MaplesRodgers2015}.
It is not a uniform microscopic estimate. Their local zero-density
moment bound, writing \(N(\sigma,t)\) for the number of zeros with
real part greater than \(\sigma\) and ordinate in \((0,t)\), is
\begin{equation}\label{eq:MR-local}
 \frac1T\int_0^T
 |N(\sigma,v+H/\log T)-N(\sigma,v)|^k\dd v
 \ll_{c,k} H^k T^{-c(\sigma-1/2)}.
\end{equation}
Here \(k\ge1\) is fixed, \(1\le H\le T^{1/4}\),
\(0<c<1/2\), and we use \(1/2<\sigma<1\)
\cite[Proposition~2.5]{MaplesRodgers2015}.
One zero with real part greater than \(\sigma\), sufficiently far
from the center-range endpoints, contributes at least
\(H/(T\log T)\) to the left-hand side for every fixed \(k\).
For fixed \(H,\sigma\), this is smaller than the permitted upper
scale: their ratio is
\(O(H^{1-k}T^{-1+c(\sigma-1/2)}/\log T)\to0\).
Thus the stated estimate allows one narrow spike. Varying \(k\)
with \(T\) would require its presently untracked constants.

Rodgers' sharp microscopic clump-frequency bound has decay
\(\exp(-cA\log A)\), averages over centers, and assumes RH
\cite[Theorem~1.1 and Corollary~1.2]{Rodgers2017Tails}.
It cannot be used as an unconditional exclusion here. These selected
source comparisons leave the inequalities in
Corollary~\ref{cor:localized-contract} unpaid; they do not assert
that no other localized theorem could help.

\subsection{Controls and the precise scope of the reduction}
The polynomial heat solution
\(H_\lambda(z)=(z-x_*)^2+2(\lambda_*-\lambda)\)
has \(g=0\), fixed midpoint and \(y^2=2s\). It attains the
\(K=0\) transport bounds and exhibits both local pair signals.
It is a local dynamical control, not a theta source; its two moving
roots do not create a broad final cluster. The smooth positive source
in Section~\ref{sec:control} also collides while satisfying the generic
analytic and operator hypotheses, but fails the literal theta arithmetic.
The repository's Euler-factor controls alter that arithmetic;
weighted-square controls lack complete reflection; fully reflected
smooth and positive-primitive controls lack the canonical square
coefficients. None is asserted to satisfy the unproved arithmetic
upper bounds proposed here. The missing discriminator must use
independent canonical information, with the same source and error
budget available to the direct scalar analysis.
